\documentclass{article}
\usepackage[T1]{fontenc}
\usepackage{lmodern}
\usepackage{graphicx}
\usepackage{amsmath,amssymb}
\usepackage[a4paper,margin=2cm]{geometry}
\usepackage[absolute,overlay]{textpos}
\setlength{\TPHorizModule}{1mm}
\setlength{\TPVertModule}{1mm}
\usepackage[indonesian]{babel}
\usepackage{hyperref}
\setlength{\emergencystretch}{2em}
% unit: Halaman 16

\begin{document}

% === Halaman 16 (llm) ===

\begin{itemize}
    \item Sebagai latihan, misalkan diberikan sebuah cipherteks:
    \begin{center}
        \texttt{TLCYKUMGDFAWTZVOYKLENSZZHYZRW}
    \end{center}
\end{itemize}

temukan kunci yang menghasilkan plainteks:
\begin{center}
    \texttt{mr johnson left his house last night}
\end{center}

lalu temukan kunci lain yang menghasilkan plainteks
\begin{center}
    \texttt{i saw the mysterious plane behind me}
\end{center}

\end{document}
